\documentclass[11pt]{article}
\usepackage[margin=1in]{geometry}
\usepackage{amsmath,amssymb,amsthm}
\newtheorem{theorem}{Theorem}
\newtheorem{lemma}[theorem]{Lemma}
\newtheorem{corollary}[theorem]{Corollary}
\newtheorem{proposition}[theorem]{Proposition}
\title{Physical factorial moments at a prime total degree}
\author{Internal proof; no priority or external-review claim}\date{17 November 2026}
\begin{document}\maketitle
Let $p$ be prime and let $n\ge1$. The integer $n$ denotes the
terminal order, so the consecutive pair has orders $n-1,n$.
Throughout this note the lower
Gamma parameter is exactly the positive integer $n$. Put
\[
 T_n(t)=\sum_{k=0}^n c_{n,k}t^k,\qquad
 c_{n,k}=\binom nk(n)_{(p-1)k}
       =\frac{n\,[n+(p-1)k-1]!}{k!(n-k)!},
 \qquad H_n(Y)=pT_n(Y/p).
\]
The second expression for $c_{n,k}$ includes $k=0$. Coefficients of
$T_n$ and $H_n$ are distinct: write $h_{n,k}=p^{1-k}c_{n,k}$ for
the latter. Normalize $v_p(p)=1$.

\begin{lemma}[Exact integral coefficients]\label{lem:coefficients}
For $k\ge1$ let
\[
 U_p(k)=\frac{(pk)!}{p^k k!}
       =\prod_{\substack{1\le j\le pk\\p\nmid j}}j,\qquad
 N=n+(p-1)k-1.
\]
Then $U_p(k)$ is a $p$-adic unit and
\[
 h_{n,k}=U_p(k)
 \left[p\binom N{pk}+\binom N{pk-1}\right].                 \tag{1}
\]
Here $\binom ab=0$ for $b>a\ge0$. Thus $H_n\in\mathbb Z[Y]$,
$h_{n,0}=p$, and $h_{n,n}=U_p(n)$ is a unit. In particular
\[
 v_p(c_{n,k})\ge k-1,\qquad v_p(c_{n,n})=n-1.              \tag{2}
\]
\end{lemma}
\begin{proof}
For $j=n-k>0$, the bracket in (1), after multiplication by
$(pk)!/(p^kk!)$, equals
\[
 \frac{N!}{p^kk!j!}\,[pj+pk]
 =\frac{pnN!}{p^kk!j!}=p^{1-k}c_{n,k}.
\]
For $j=0$ its first binomial is zero and its second binomial is
one, giving the same identity. The product formula for $U_p(k)$
proves all remaining assertions. No factorial valuation estimate
or analytic moment-generating series is used.
\end{proof}

\begin{theorem}[An Eisenstein index family]\label{thm:index}
If $n=bp^r$ with $r\ge0$ and $1\le b<p$, then $H_n$ is Eisenstein
at $p$. Consequently $T_n$ is irreducible over $\mathbb Q_p$
and over $\mathbb Q$.
\end{theorem}
\begin{proof}
The constant coefficient has valuation one and the leading
coefficient is a unit by Lemma~\ref{lem:coefficients}. For
$0<k<n$ it remains to show
\[
 p\mid \binom{n+(p-1)k-1}{pk-1}.                          \tag{3}
\]
Recall the elementary no-carry form of Lucas' criterion:
$\binom{a+b}{a}$ is a $p$-adic unit precisely when adding $a$
and $b$ in base $p$ produces no carry. This follows by taking
coefficients in
$(1+X)^m=\prod_i(1+X^{p^i})^{m_i}$ modulo $p$.
If (3) failed, adding $pk-1$ and $n-k$ would be carry-free.
Put $s=v_p(k)$. The last $s+1$ digits of $pk-1$ are $p-1$,
so $p^{s+1}\mid n-k$. Therefore $v_p(n)=s$, hence $s=r$.
But then $k=ap^r$ with $1\le a<b<p$, and
$n-k=(b-a)p^r$ is not divisible by $p^{r+1}$, a contradiction.
Equation (1) now makes every interior coefficient $p$-divisible.
The Eisenstein argument is over the DVR $\mathbb Z_p$ and
allows a unit leading coefficient. The invertible variable change
$t=Y/p$ preserves irreducibility over $\mathbb Q_p$.
\end{proof}

\begin{corollary}[Ordinary consecutive moments]\label{cor:pair}
For every prime $p$ and every terminal index $n=bp^r\ge2$ with
$1\le b<p$, the ordinary factorial moments of every nonzero phase
\[
 f(X)=X(A+BX^{p-1})
\]
at orders $n-1,n$ cannot both vanish. In particular this proves
the gap-six assertion at all terminal $n=b7^r$,
$1\le b\le6$, $n\ge2$.
\end{corollary}
\begin{proof}
If $A B=0$ this is immediate for a monomial. Otherwise put
$t=B/A$ and $\theta=tD_t$. The normalized upper and lower
moments, respectively, are
\[
 T_n+\frac{p-1}{n}\theta T_n,\qquad
 T_n-\frac1n\theta T_n.
\]
These identities follow coefficientwise from
$(n+1)_m=(1+m/n)(n)_m$ and
$\binom{n-1}{k}=(1-k/n)\binom nk$.
The $2$ by $2$ coefficient determinant is $-p/n\ne0$.
Since $T_n(0)=1$, a common zero would be a multiple root of
$T_n$. Theorem~\ref{thm:index} excludes this in characteristic zero.
\end{proof}

\begin{theorem}[The exact possible critical radius]\label{thm:radius}
For arbitrary $n\ge1$, every positive-valuation root of $H_n$
is simple. More precisely, all positive-valuation roots form one
irreducible Eisenstein factor over $\mathbb Q_p$.
Thus every multiple root $t$ of $T_n$ would satisfy
\[
 v_p(pt)=0.
\]
Equivalently, for $\epsilon=(-1)^p$ and $u=\epsilon/t$,
\[
                         v_p(u)=1.                       \tag{4}
\]
This is necessary, not sufficient, for a multiple root.
\end{theorem}
\begin{proof}
Divide $H_n$ by its unit leading coefficient. It is monic
integral with constant valuation one, so all its roots are
integral. Write
$\overline H_n=Y^m G(Y)$, where $m\ge1$ and $G(0)\ne0$.
The two coprime factors lift by Hensel's lemma to monic integral
factors $H_+,H_0$, with $\overline H_+=Y^m$ and
$\overline H_0=G$ (up to the harmless unit normalization).
Every root of $H_+$ has positive valuation, every root of $H_0$
has valuation zero, and $H_0(0)$ is a unit. Hence $H_+(0)$ has
valuation one. Thus $H_+$ is Eisenstein and irreducible and,
in characteristic zero, has only simple roots.
A multiple root of $H_n$ must therefore have valuation zero.
The change $Y=pt$ and then $u=\epsilon/t$ gives (4).
Alternatively, a monic irreducible positive-valuation factor
occurring twice would contribute at least two to the constant
valuation, contradicting its value one.
\end{proof}

\begin{proposition}[A finite-field trace reduction]\label{prop:trace}
Coefficientwise modulo $p$,
\[
 \overline H_n(Y)=\sum_{k=1}^n(-1)^k
        \binom{n+(p-1)k-1}{pk-1}Y^k,\qquad
 \sum_{n\ge1}\overline H_n(Y)z^n
       =\frac{-Yz}{1+Yz-z^p}.                             \tag{5}
\]
In particular every nonzero exponent is congruent to $n$
modulo $p$. For odd $p$, put $w=(p-1)!$ and define $K_n$ by
\[
 K_n(Y)=\sum_{k=0}^n a_{n,k}Y^k,\quad a_{n,0}=p,\quad
 a_{n,k}=p\binom{n+(p-1)k-1}{pk}
              +\binom{n+(p-1)k-1}{pk-1}\quad(k>0).
\]
Then
\[
 H_n(Y)\equiv K_n(wY)\pmod {p^2}.                         \tag{6}
\]
If $D(z)=(1-z)^p-Yz$, the exact generating function is
\[
 \sum_{n\ge0}K_n(Y)z^n=p-\frac{zD'(z)}{D(z)},\quad K_0=p. \tag{7}
\]
Thus $K_n(Y)$ is the sum of the $n$th powers of the $p$ roots,
with multiplicity, of $(\lambda-1)^p-Y\lambda^{p-1}$.
\end{proposition}
\begin{proof}
Wilson's congruence gives $U_p(k)\equiv(-1)^k\pmod p$.
The unit test for the binomial in (5) forces $p\mid n-k$,
since $pk-1$ ends in the digit $p-1$. For its generating function,
write $n=k+j$ and use
$\sum_{j\ge0}\binom{pk+j-1}{pk-1}z^j=(1-z)^{-pk}$,
then sum the geometric series in $-Yz/(1-z)^p$.

For odd $p$, each block of nonmultiples of $p$ satisfies
\[
 \prod_{i=1}^{p-1}(p\ell+i)
 \equiv(p-1)!\left(1+p\ell\sum_{i=1}^{p-1}i^{-1}\right)
 \equiv(p-1)!\pmod {p^2}.
\]
The inverse sum is zero modulo $p$. Therefore
$U_p(k)\equiv w^k\pmod {p^2}$, proving (6).
For (7), put $Q=Yz/(1-z)^p$. The two binomial
sums and the constant coefficients give exactly
\[
 \sum_{n\ge0}K_n(Y)z^n
   =\frac{p}{1-z}
      +\left(\frac{pz}{1-z}+1\right)\frac{Q}{1-Q}.
\]
Simplifying yields
\[
 \frac{p(1-z)^{p-1}-(p-1)Yz}{(1-z)^p-Yz}
       =p-\frac{zD'(z)}{D(z)}.
\]
The power-sum interpretation follows by factoring $D$ formally
over an algebraic splitting field.
\end{proof}

\section*{Two precise limitations of the residue reduction}
\begin{proposition}[Unit-residue obstructions]\label{prop:local-obstructions}
Neither the unique positive-valuation factor nor its odd/even
degree controls the unit factors in Theorem~\ref{thm:radius}.
For example, at $p=7,n=15$,
\[
 \overline H_{15}=6Y(Y-3)^7(Y-5)^7.
\]
For the $Y=3+Z$ residue disk the coefficients through degree seven
have residues modulo $7^3=343$
\[
 [49,210,182,14,70,280,273,153].
\]
Their valuations are $[2,1,1,1,1,1,1,0]$.
The local degree-seven factor consequently has Newton edges
$(0,2)$ to $(1,1)$ to $(7,0)$. Its local factors have degrees
one and six, respectively; the latter is irreducible because
its slope denominator is six. Thus a universal odd-degree local
factor assertion is false, even on the physical integer family.
This is not an SFC counterexample.

Even first-order transversality after division by $p$ is not
uniform. At $p=7,n=145$, the first two coefficients of
$H_{145}(4+Z)$ modulo $7^4=2401$ are $1470,1176$.
Both have valuation exactly two:
\[
 H_{145}(4)\equiv0,\qquad H_{145}'(4)\equiv0\pmod {49}.
\]
The coefficients of $Z^0,\ldots,Z^{14}$ in
$\overline H_{145}(4+Z)$ are
$[0,0,0,0,0,0,0,0,0,0,0,0,0,0,5]$,
so the unit residue root at $4$ has multiplicity fourteen. Thus neither a uniform simple residue root
nor the universal criterion
$v_7(H_n(a))\ge2\Rightarrow v_7(H'_n(a))=1$
can prove all indices. The actual polynomial still has no
certified multiple root here; these are obstructions to shortcuts,
not counterexamples to simplicity.

\end{proposition}
\begin{proof}
The displayed coefficients follow from the exact integer identity (1),
using the finite binomial expansion of $H_n(a+Z)$. Their
valuations give the stated Newton edges; the slope denominator
forces the degree-six factor to be irreducible. The last residue
list proves multiplicity fourteen at $4$. These exact finite
calculations are replayed by the accompanying checker.
\end{proof}

The ordinary gap-six all-index assertion remains open.
The uniform reductions above retain the exact physical
parameter, locate its remaining multiple-root problem at the
single radius (4), and give an unbounded Eisenstein index family.
\end{document}
